\section*{Chapter \ref{ch:g12:ph-conductivity-titrations} --- Titrations by pH and Conductivity}

\begin{solution}{exo:g12:ph-conductivity-titrations:1}
$V_E = \qty{20.0}{mL}$; $c = 0.100 \times 20.0 / 20.0 = \qty{0.100}{mol/L}$.
\end{solution}

\begin{solution}{exo:g12:ph-conductivity-titrations:2}
At \omterm{def:g12:ph-conductivity-titrations:ph-metric}{half-equivalence} \omterm{def:g9:acids-bases-ph:ph}{pH} $= pK_a$: $pK_a = 3.75$, close to methanoic acid
(3.74).
\end{solution}

\begin{solution}{exo:g12:ph-conductivity-titrations:3}
Phenolphthalein: its range 8.0--10.0 lies in the jump, which is basic.
Methyl red (4.2--6.3) would change near \omterm{def:g12:ph-conductivity-titrations:ph-metric}{half-equivalence}, far too early.
\end{solution}

\begin{solution}{exo:g12:ph-conductivity-titrations:4}
The probe's response drifts with time and temperature: it is set to read
correctly in two \omterm{def:g12:buffers-predominance:buffer}{buffer solutions} of known \omterm{def:g9:acids-bases-ph:ph}{pH}.
\end{solution}

\begin{solution}{exo:g12:ph-conductivity-titrations:5}
\ce{H3O+} (349.8) > \ce{OH-} (198.3) > \ce{Cl-} (76.2) > \ce{Na+} (50.3).
\end{solution}

\begin{solution}{exo:g12:ph-conductivity-titrations:6}
$c = 0.050 \times 14.6 / 20.0 = \qty{0.0365}{mol/L}$.
\end{solution}

\begin{solution}{exo:g12:ph-conductivity-titrations:7}
Slopes: 2.0, 4.5, 17.5 and \qty{2.5}{mL^{-1}}. The largest lies between
10.0 and \qty{10.2}{mL}: $V_E \approx \qty{10.1}{mL}$.
\end{solution}

\begin{solution}{exo:g12:ph-conductivity-titrations:8}
Before \omterm{def:g11:titration:equivalence}{equivalence} each \ce{OH-} added removes an \ce{H3O+} (very
conducting) and brings a \ce{Na+} (much less conducting): the
\omterm{def:g12:ph-conductivity-titrations:conductivity}{conductivity} falls. After it, \ce{Na+} and \ce{OH-} accumulate without
reacting: it rises.
\end{solution}

\begin{solution}{exo:g12:ph-conductivity-titrations:9}
At the start, few \omterm{def:g9:ions:ion}{ions} (the \omterm{def:g12:ka-and-pka:strong-acid}{weak acid} hardly reacts with water): low
\omterm{def:g12:ph-conductivity-titrations:conductivity}{conductivity}. Before \omterm{def:g11:titration:equivalence}{equivalence}, ethanoate and sodium \omterm{def:g9:ions:ion}{ions} are formed
in place of uncharged \omterm{def:g7:atoms-and-molecules:molecule}{molecules}: it rises, slowly. After \omterm{def:g11:titration:equivalence}{equivalence},
\ce{Na+} and \ce{OH-} accumulate: it rises faster. Two rising lines,
the second steeper.
\end{solution}

\begin{solution}{exo:g12:ph-conductivity-titrations:10}
About 2.9, 4.76 and 8.7. Ethanoic acid is weak: only a few percent of it
gives \omterm{def:g12:ka-and-pka:oxonium}{oxonium ions}, so the \omterm{def:g9:acids-bases-ph:ph}{pH} starts higher than the 1.0 of hydrochloric
acid at a similar concentration.
\end{solution}

\begin{solution}{exo:g12:ph-conductivity-titrations:11}
So that the volume of \omterm{def:g11:titration:titration}{titrant} added hardly changes the total volume: the
concentrations, and so the \omterm{def:g12:ph-conductivity-titrations:conductivity}{conductivity}, then vary along straight lines.
\end{solution}

\begin{solution}{exo:g12:ph-conductivity-titrations:12}
The hydroxide \omterm{def:g9:ions:ion}{ions} react first with the \omterm{def:g12:ka-and-pka:strong-acid}{strong acid} (\ce{H3O+}), then
with the weak one. The first jump marks the end of the hydrochloric acid
(its \omterm{def:g11:titration:equivalence}{equivalent volume} measures it), the second the end of the ethanoic
acid (the volume between the two jumps measures it).
\end{solution}

\begin{solution}{exo:g12:ph-conductivity-titrations:13}
Start: $(349.8 + 76.2) \times 0.0100 = \qty{4.26}{mS/cm}$. At
\omterm{def:g11:titration:equivalence}{equivalence}: $[\ce{Na+}] = [\ce{Cl-}] = 1.00/110 = \qty{0.00909}{mol/L}$,
so $(50.3 + 76.2) \times 0.00909 = \qty{1.15}{mS/cm}$. Both as on the
figure.
\end{solution}

\begin{solution}{exo:g12:ph-conductivity-titrations:14}
None on $V_E$: the jump is at the same volume, whatever the offset of the
readings. The $pK_a$ read at \omterm{def:g12:ph-conductivity-titrations:ph-metric}{half-equivalence} is 0.2 too high.
\end{solution}

\begin{solution}{exo:g12:ph-conductivity-titrations:15}
The \omterm{def:g9:acids-bases-ph:ph}{pH} does not change during this \omterm{def:g9:ions:precipitate}{precipitation}, but the \omterm{def:g9:ions:ion}{ions} do:
before \omterm{def:g11:titration:equivalence}{equivalence} each \ce{Ag+} added removes a \ce{Cl-} and brings a
\ce{NO3-} of similar \omterm{def:g12:ph-conductivity-titrations:conductivity}{conductivity}, so the \omterm{def:g12:ph-conductivity-titrations:conductivity}{conductivity} stays nearly
constant (the sodium \omterm{def:g9:ions:ion}{ions} were there from the start); after \omterm{def:g11:titration:equivalence}{equivalence} \ce{Ag+} and \ce{NO3-}
accumulate and it rises. A flat line, then a rising one.
\end{solution}

\begin{solution}{pb:g12:ph-conductivity-titrations:1}
\textbf{1.} \qty{6}{g} of ethanoic acid per \qty{100}{g} of vinegar.

\textbf{2.} \qty{101}{g} of vinegar hold $6.0 \times 101/100 =
\qty{6.06}{g}$; $M = \qty{60.0}{g/mol}$: \qty{0.101}{mol} in
\qty{0.1000}{L}, that is \qty{1.01}{mol/L}.

\textbf{3.} Undiluted, \qty{10.0}{mL} would need about \qty{101}{mL} of
sodium hydroxide at \qty{0.100}{mol/L}: more than a burette holds.

\textbf{4.} A \qty{10.0}{mL} volumetric pipette and a \qty{100.0}{mL}
volumetric flask.

\textbf{5.} \ce{CH3COOH + OH- -> CH3COO- + H2O}.

\textbf{6.} $V_E = \qty{10.1}{mL}$.

\textbf{7.} $c = 0.100 \times 10.1 / 10.0 = \qty{0.101}{mol/L}$.

\textbf{8.} $10 \times 0.101 = \qty{1.01}{mol/L}$.

\textbf{9.} \omterm{def:g9:acids-bases-ph:ph}{pH} 4.76 at \qty{5.05}{mL}: the $pK_a$ of ethanoic acid.

\textbf{10.} At \omterm{def:g11:titration:equivalence}{equivalence} the solution holds sodium ethanoate, and
ethanoate is a \omterm{def:g12:ka-and-pka:strong-acid}{weak base}.

\textbf{11.} Phenolphthalein, whose range 8.0--10.0 contains the \omterm{def:g9:acids-bases-ph:ph}{pH} at
\omterm{def:g11:titration:equivalence}{equivalence} (about 8.7).

\textbf{12.} Sodium \omterm{def:g9:ions:ion}{ions} and ethanoate \omterm{def:g9:ions:ion}{ions}.

\textbf{13.} Each \ce{OH-} added turns a nearly un-ionised \omterm{def:g7:atoms-and-molecules:molecule}{molecule} into
\ce{CH3COO-}, with a \ce{Na+}: \omterm{def:g9:ions:ion}{ions} of modest \omterm{def:g12:ph-conductivity-titrations:conductivity}{conductivity} are added
slowly.

\textbf{14.} Beyond \omterm{def:g11:titration:equivalence}{equivalence} the hydroxide \omterm{def:g9:ions:ion}{ions}, very conducting,
accumulate with the sodium \omterm{def:g9:ions:ion}{ions}.

\textbf{15.} The \omterm{def:g12:ka-and-pka:strong-acid}{weak acid} gives very few \omterm{def:g9:ions:ion}{ions} to water.

\textbf{16.} At $V_E = \qty{10.1}{mL}$, as in the \omterm{def:g12:ph-conductivity-titrations:ph-metric}{pH-metric titration}.

\textbf{17.} $1.01 \times 60.0 = \qty{60.6}{g/L}$.

\textbf{18.} \qty{6.06}{g} in \qty{100.0}{mL}.

\textbf{19.} $6.06 \times 100 / 101 = \qty{6.0}{g}$ per \qty{100}{g}.

\textbf{20.} The vinegar has an acidity of 6.0 degrees: the label is
right.
\end{solution}
